\section{Игра в билеты}\label{sec:tickets}

Каждый из предыдущих разделов рассматривает одну сторону поворота. Этот раздел собирает их
в одной комнате. Всё в нём элементарно, и каждое число проверено скриптом \texttt{verify.py}.
Смысл не в новой математике, а в том, что три пути, энергия, близнецы и треугольник
Бертрана появляются, стоит только задать простой вопрос о ходьбе.

\paragraph{Правила.} Все начинают в центре $O$ комнаты. Каждому выдают
билеты, и билет оплачивает один ход по прямой длины не более $1$. Два игрока, $A$ и
$B$, тратят по одному билету, уходя из $O$ каждый в выбранном им направлении, так что
они оказываются на единичной окружности под углом $\theta$ друг к другу. Теперь $A$ хочет снова встретиться с $B$. Одно
правило не ослабляется никогда: ни один ход не может стоить больше одного билета.

\paragraph{Три цены одного перехода.} У перехода есть три естественные цены, по одной
на каждый способ его измерить.
\begin{itemize}
  \item \emph{Дуга}, путь, которым шёл поворот, стоит $\theta$ билетов. Она укладывается в один билет
    вплоть до ровно одного радиана, $57.30^\circ$. Это и есть радиан: угол, дуга которого
    имеет длину в один билет.
  \item \emph{Хорда}, напрямик, стоит $2\sin(\theta/2)$. Она укладывается в один билет вплоть
    до ровно $60^\circ$. Тогда $A$, $B$ и $O$ образуют равносторонний треугольник, а шесть
    игроков, расставленных через $60^\circ$, стоят в вершинах правильного шестиугольника, стороны которого
    равны его радиусу. С этого шестиугольника начинал Архимед.
  \item \emph{Энергия} в решёточном счёте "--- это энергия грани Вильсона
    $1-\cos\theta=\tfrac12C^2$. В этих единицах один ход длины $L$ стоит
    $\tfrac12L^2$, так что два хода по радиусам обратно через центр вместе стоят
    $\tfrac12+\tfrac12=1$. Назовём это одним билетом энергии. Переход стоит
    ровно один билет энергии при $90^\circ$, меньше при меньших углах и больше при больших, вплоть до двух на
    противоположной стороне комнаты.
\end{itemize}
У каждой цены свой угол в один билет, и все три различны.
\begin{center}\small
\begin{tabular}{llll}\toprule
Цена & Стоимость перехода & Один билет при & Почему \\ \midrule
дуга & $\theta$ & $57.30^\circ$ & определение радиана \\
хорда & $2\sin(\theta/2)$ & $60^\circ$ & равносторонний треугольник \\
энергия & $1-\cos\theta$ & $90^\circ$ & Пифагор \\
\bottomrule\end{tabular}
\end{center}

\paragraph{Вся доска.} Таблица выше оценивает переход тремя способами, но у каждого
пути есть и длина, и энергия, так что на доске шесть клеток, и каждая однозначно определена,
как только зафиксировано её правило. Если энергию измерять как $\tfrac12L^2$, как выше:
\begin{center}\small
\begin{tabular}{lccc}\toprule
Один билет\ldots & внешний путь & дуга & хорда \\ \midrule
длины & $53.13^\circ$ & $57.30^\circ$ & $60^\circ$ \\
 & треугольник $3$--$4$--$5$ & один радиан & шестиугольник \\[2pt]
энергии & $70.53^\circ$ & $81.03^\circ$ & $90^\circ$ \\
 & $\cos\theta=\tfrac13$ & $\sqrt2$ радиан & квадрат \\
\bottomrule\end{tabular}
\end{center}
В каждой строке внешний путь, самый дорогой, покупает наименьший поворот, а хорда, самая
дешёвая, "--- наибольший. Две клетки сами по себе являются классическими углами.
Внешний путь покупает один билет длины, когда его отрезок касательной $\tan(\theta/2)=\tfrac12$, и
построение Кэли с таким отрезком даёт точку $(\tfrac35,\tfrac45)$. При любом рациональном
отрезке $t$ построение даёт рациональную точку $\bigl((1-t^2)/(1+t^2),\,2t/(1+t^2)\bigr)$,
и именно так возникает каждая пифагорова тройка. Внешний путь покупает один билет энергии при
$\cos\theta=\tfrac13$ "--- это двугранный угол правильного тетраэдра.

\paragraph{Зазор.} Между $57.30^\circ$ и $60^\circ$ короткий путь напрямик укладывается в один
билет, а истинный путь "--- нет. Величина этого зазора при угле шестиугольника равна
$\pi/3-1\approx0.0472$ билета, надбавка в $4.72\%$, и это то же число, что
отношение длины окружности к периметру шестиугольника, $2\pi/6$, минус один. Это
цена кривизны в масштабе одного билета. Именно это вычисляется в разделе~\ref{sec:lengths},
$A-C\approx\theta^3/24$, только увиденное в размере, который можно пройти ногами.

\paragraph{Почему энергия окупается именно при $90^\circ$.} Квадрат перехода равен
$|A-B|^2=|A|^2+|B|^2-2A\cdot B=2-2\cos\theta$. Путь обратно через центр состоит из двух
отрезков длины один. По энергии переход и два отрезка стоят одинаково ровно тогда, когда
$A\cdot B=0$: при прямом угле, где хорда "--- гипотенуза двух билетов.
Ниже $90^\circ$ срезать напрямик выгодно по энергии. Выше $90^\circ$ два отдельных хода
через центр обходятся по энергии дешевле, чем один переход. Энергия растёт как
квадрат длины, поэтому разбиение длинного хода на два коротких экономит энергию, и за
прямым углом разбиение выигрывает.

\paragraph{Целые билеты: почему всегда два хода.} Теперь ужесточим правило: билет
нужно тратить целиком, так что каждый ход имеет длину ровно один. С одним целым билетом $A$
может попасть только на окружность радиуса один вокруг $A$. Если $B$ не стоит точно там
(при $\theta=60^\circ$), $A$ перелетает или недотягивает. С двумя целыми билетами $A$ может
достичь любой точки на расстоянии не больше двух, но вершина, в которой $A$ поворачивает, должна лежать в одном билете
от $A$ и в одном билете от $B$. Две окружности радиуса один вокруг $A$ и $B$ пересекаются ровно
в двух точках, симметричных относительно хорды. Одна из них "--- центр $O$; другая "--- его
отражение относительно середины хорды, точка $A+B$. Итак, существует ровно два
маршрута из двух билетов:
\begin{itemize}
  \item \emph{через центр}: дойти домой, затем выйти по пути $B$;
  \item \emph{через вершину-близнеца} $A+B$, вершину, противоположную $O$ в ромбе
    $O$, $A$, $A+B$, $B$, все четыре стороны которого имеют длину один.
\end{itemize}
Эти две вершины "--- близнецы в смысле раздела~\ref{sec:twins}: два ответа на один
вопрос, зеркальные отражения друг друга относительно хорды. При полуобороте они сливаются, поскольку $A+B=O$
при $B=-A$, и остаётся только маршрут через центр.

\paragraph{Где стоит вершина-близнец.} Вершина-близнец находится на расстоянии
$|A+B|=2\cos(\theta/2)$ от центра, вдвое дальше середины хорды. Она
лежит вне круга при углах меньше $120^\circ$, внутри при больших и ровно на окружности при
$120^\circ$. При этом угле переход равен $\sqrt3$, стороне треугольника Бертрана, и
вершина-близнец попадает в вершину шестиугольника между двумя игроками: ромб распадается
на два равносторонних треугольника.

\paragraph{Где стоит вершина обхода.} Внешний путь из раздела~\ref{sec:picture}
поворачивает в вершине касательных $Q$, на расстоянии $\sec(\theta/2)$ от центра на той же
биссектрисе. Вершина касательных и вершина-близнец совпадают ровно тогда, когда
$\sec(\theta/2)=2\cos(\theta/2)$, то есть при $\theta=90^\circ$. Тогда ромб "--- это
квадрат, его дальняя вершина "--- точка пересечения касательных, каждый отрезок касательной равен $\tan45^\circ=1$,
и внешний путь составляет ровно два целых билета. Ниже $90^\circ$ вершина касательных "---
ближняя из двух, и обход стоит меньше двух билетов; выше $90^\circ$ она
дальняя, и обход стоит больше.

\paragraph{Близнец прямого угла.} Энергия перехода к $B$ равна $1-\cos\theta$;
энергия перехода к точке, противоположной $B$, равна $1+\cos\theta$. В сумме они всегда дают два
билета, потому что $A$, $B$ и $-B$ образуют прямой угол при $A$ (Фалес). В обозначениях
раздела~\ref{sec:number} эти две энергии равны в точности $2d^2$ и $2S$, так что это снова $S+d^2=1$,
а точка безубыточности при $90^\circ$ "--- это точка, где $S=d^2=\tfrac12$: два
положения ровно настолько же одинаковы, насколько различны.

\paragraph{Какое правило энергии.} Прямой угол "--- точка безубыточности для правила Вильсона, энергия
$\tfrac12C^2$. Это правило, а не факт о комнате. С правилом континуума, энергия
$\tfrac12\theta^2$, точка безубыточности сдвигается к $81.03^\circ$; с правилом внешнего пути,
$\tfrac12T^2$, "--- к $70.53^\circ$. Выбор между ними определяется пределом малых билетов, где
правила обязаны согласовываться с континуумом: при $5^\circ$ три энергии равны $0.003805$,
$0.003808$ и $0.003813$, в пределах $0.2\%$, и они становятся одинаковыми по мере того, как билеты уменьшаются.
Именно на больших поворотах решёточные правила расходятся, и раздел~\ref{sec:energy} посвящён как раз
этому различию.

Итак, прямой угол особенный трижды. По правилу Вильсона переход стоит один
билет энергии, и его энергия равна энергии двух ходов через центр. Вообще без всякого правила,
из одной лишь геометрии, обход становится двумя целыми билетами, через вершину квадрата.

\begin{table}[h!]\centering\small
\caption{Игра в билеты при пяти углах: три цены перехода, его энергия
Вильсона и расстояния от центра до вершины-близнеца $A+B$ и до вершины
касательных $Q$.}
\label{tab:tickets}
\begin{tabular}{rrrrrrr}\toprule
$\theta$ & дуга & хорда & внешний путь & энергия & вершина-близнец & вершина касательных \\ \midrule
$57.30^\circ$ & 1.0000 & 0.9589 & 1.0926 & 0.4597 & 1.7552 & 1.1395 \\
$60^\circ$ & 1.0472 & 1.0000 & 1.1547 & 0.5000 & 1.7321 & 1.1547 \\
$90^\circ$ & 1.5708 & 1.4142 & 2.0000 & 1.0000 & 1.4142 & 1.4142 \\
$120^\circ$ & 2.0944 & 1.7321 & 3.4641 & 1.5000 & 1.0000 & 2.0000 \\
$180^\circ$ & 3.1416 & 2.0000 & $\infty$ & 2.0000 & 0 & $\infty$ \\
\bottomrule\end{tabular}
\end{table}

\begin{figure}[h!]\centering
\begin{tikzpicture}[scale=2.6]
  \def\t{70}
  \coordinate (O) at (0,0);
  \coordinate (A) at (1,0);
  \coordinate (B) at ({cos(\t)},{sin(\t)});
  \coordinate (P) at ({1+cos(\t)},{sin(\t)});
  \coordinate (Q) at (1,{tan(\t/2)});
  \draw[gray!50] (O) circle (1);
  \draw[thick] (O) -- (A); \draw[thick] (O) -- (B);
  \draw[gray,dotted,thick] (A) -- (P) -- (B);
  \draw[very thick,dgarc] (A) arc (0:\t:1);
  \draw[very thick,dgchord] (A) -- (B);
  \draw[very thick,dgtan] (A) -- (Q) -- (B);
  \fill (O) circle (0.022) node[below left] {$O$};
  \fill (A) circle (0.022) node[below right] {$A$};
  \fill (B) circle (0.022) node[above left] {$B$};
  \fill[gray] (P) circle (0.022) node[above right] {$A+B$};
  \fill[dgtan] (Q) circle (0.022) node[right] {$Q$};
  \draw (0.2,0) arc (0:\t:0.2); \node at ({0.29*cos(\t/2)},{0.29*sin(\t/2)}) {$\theta$};
\end{tikzpicture}
\caption{Игра в билеты при $\theta=70^\circ$. Каждый игрок прошёл от
центра на один билет. Хорда (синяя), дуга (зелёная) и внешний путь через вершину касательных $Q$
(оранжевый) "--- три цены перехода. Два маршрута из целых билетов поворачивают в
центре или в его близнеце $A+B$, дальней вершине единичного ромба (пунктир). При $90^\circ$
вершина касательных и вершина-близнец совпадают.}
\label{fig:tickets}
\end{figure}

\begin{plain}
Представьте комнату. Все начинают в середине с билетами, и один билет оплачивает один
проход по прямой не длиннее одного шага. Двое друзей тратят по билету и расходятся в
разные стороны. Как одному вернуться к другому? Пройти по кривой стоит один
билет, только если между ними не больше примерно $57^\circ$. Срезать напрямик стоит один
билет, если между ними не больше $60^\circ$. Если считать так, как физики считают энергию,
прямой переход стоит один билет ровно при $90^\circ$ "--- столько же, сколько вернуться в
середину и снова выйти; это Пифагор. Если билеты нужно тратить целиком, всегда есть
ровно два способа сделать это за два хода: повернуть в середине комнаты или повернуть в
её зеркальном углу по другую сторону перехода. Эти два угла "--- близнецы. При
$120^\circ$ угол-близнец попадает прямо на окружность, а при $90^\circ$ он становится
вершиной идеального квадрата, где встречаются две боковые линии и внешний путь равен ровно
двум целым билетам.
\end{plain}
